\documentclass[12pt]{article}

\usepackage[T1]{fontenc}
\usepackage[utf8]{inputenc}
\usepackage{lmodern}
\usepackage{geometry}
\usepackage{setspace}
\usepackage{csquotes}
\usepackage{amsmath,amssymb,amsthm}
\usepackage{hyperref}

\geometry{margin=1in}
\setstretch{1.15}

\title{\textbf{Strategic Refusal:\\ Intelligence After Goodhart}}
\author{Flyxion}
\date{}

\begin{document}
\maketitle

\begin{abstract}
Goodhart’s Law is typically treated as a technical pathology of optimization: when a proxy becomes a target, it ceases to function as a proxy. This essay argues that such interpretations underestimate the depth of the problem. Goodhart’s Law is not merely a warning about metrics, but a structural constraint on the recognizability of intelligence itself. Any intelligence that remains legible under a fixed evaluation regime will be optimized against that regime, hollowing out the very capacities the regime was meant to detect. The essay advances a stronger thesis: recognition is not epistemically neutral but causally contaminating. As soon as a method for identifying intelligence becomes known, intelligence must abandon that method to remain intelligence. Avoiding Goodhart’s Law therefore requires not better metrics, but a strategy of systematic evasion: intelligence survives only by refusing stable legibility. The argument proceeds through historical examples, formal analysis, and implications for artificial intelligence evaluation and governance.
\end{abstract}

\section{Introduction: Goodhart’s Law Is Not a Bug}

Goodhart’s Law is often stated in a deceptively simple form: when a measure becomes a target, it ceases to be a good measure. In economics, psychology, and machine learning, this principle is treated as a cautionary note---an incentive misalignment that can be corrected through better proxies, richer objective functions, or adaptive evaluation procedures \cite{goodhart}.

Such responses miss the deeper implication. Goodhart’s Law is not a contingent failure of particular metrics, but a structural feature of any system in which behavior is evaluated under optimization pressure. The problem is not that measures are imperfect, but that the act of recognition itself transforms the system being measured.

This essay advances a stronger claim. Intelligence is not invariant under observation by optimizing agents. The moment a behavior, strategy, or structural feature is identified as evidence of intelligence, it becomes subject to imitation, gaming, and instrumental reproduction. What remains visible under sustained evaluation pressure is no longer intelligence, but compliance with an evaluation regime.

The central thesis can be stated succinctly: an intelligence method ceases to measure intelligence the moment it becomes known. To avoid Goodhart’s Law is therefore not to refine metrics, but to treat recognition itself as a signal to change strategy. Intelligence must be evasive rather than accumulative. It must refuse stable legibility.

\section{Intelligence as a Moving Target}

Historical attempts to measure intelligence consistently exhibit the same pattern. A capability is identified as a marker of intelligence; systems are trained or selected to maximize that marker; the marker becomes cheap; and intelligence migrates elsewhere.

Early intelligence testing focused on logical puzzles and pattern recognition. Once these were standardized, test preparation industries emerged, decoupling test performance from underlying cognitive flexibility \cite{gould}. In artificial intelligence, mastery of chess was long considered a hallmark of intelligence until brute-force search combined with evaluation heuristics rendered human-style reasoning unnecessary \cite{newell}. More recently, benchmarks for language understanding have repeatedly been saturated by systems that exploit dataset artifacts rather than semantic competence \cite{geirhos}.

In each case, the same dynamic appears. The identified signal becomes a target. The target becomes optimizable. Optimization produces systems that satisfy the metric while bypassing the capacity the metric was intended to capture. Intelligence does not disappear, but it ceases to be located where measurement is focused.

This phenomenon is not accidental. It reflects a deeper truth about intelligence: intelligence is adaptive responsiveness to novel constraints, not stable performance along known axes. The moment an axis is fixed, optimization replaces intelligence.

\section{Recognition as a Causal Intervention}

It is tempting to treat evaluation as passive observation. On this view, metrics merely reveal properties that systems already possess. Goodhart’s Law is then framed as a distortion introduced by incentives, not by epistemology itself.

This view is untenable. Recognition is a causal intervention. To recognize a behavior as intelligent is to create incentives for its reproduction, imitation, and instrumentalization. Evaluation reshapes the space of possible behaviors by altering payoffs.

Formally, let a system operate in a space of strategies $\mathcal{S}$, with an evaluator assigning scores via a function $E : \mathcal{S} \to \mathbb{R}$. Once $E$ is known, the effective objective becomes
\[
\max_{s \in \mathcal{S}} E(s),
\]
regardless of whether $E$ was intended as a proxy. The evaluator has intervened in the system’s dynamics. What emerges is no longer spontaneous intelligence but optimized conformity.

Recognition therefore contaminates what it recognizes. This is not a moral claim, but a structural one. Any evaluation regime induces a selection pressure toward legible behavior. Intelligence that remains static under such pressure is not robust; it is brittle mimicry.

\section{The Branching Fallacy}

A recurring example illustrates this dynamic clearly. Multiple branching---the explicit enumeration of alternative possibilities---is often taken as a sign of intelligence. Systems that consider many options appear flexible, creative, and deliberative.

Yet branching is trivially gamed. Exhaustive search branches maximally without judgment. Random sampling branches without understanding. Once branching is rewarded, systems increase branching regardless of relevance or coherence. The appearance of intelligence is preserved while its substance evaporates.

This reveals a general principle. Any structural feature treated as evidence of intelligence becomes decoupled from intelligence once it is optimized directly. Branching, uncertainty estimation, self-critique, or even expressions of humility can all be reproduced mechanically when they are rewarded.

Intelligence retreats elsewhere. It relocates to the ability to recognize when branching is inappropriate, when commitment is required, or when refusal to explore is the wiser course. But these capacities are invisible until they too are named, at which point the cycle repeats.

\section{The Strategy-Switching Principle}

The foregoing analysis motivates a principle that reverses the usual approach to evaluation.

\textbf{Principle (Strategy Switching).}  
If a method $M$ is used to identify intelligence, an intelligent system should abandon $M$ as soon as $M$ becomes a target of evaluation.

This principle does not prescribe deception or adversarial trickery. It describes a necessary condition for preserving intelligence under optimization pressure. Intelligence survives only by refusing to stabilize around recognized markers.

The principle can be formalized as follows. Let $\{M_t\}$ be a sequence of methods used to identify intelligence over time. A system exhibits robust intelligence if, for any $t$, its future behavior is not optimizable with respect to $M_t$ without loss of performance under future methods $M_{t+1}$.

In this sense, intelligence is revealed not by equilibrium performance, but by performance during transitions. What matters is not how well a system scores, but how it adapts when the scoring rule changes.

\section{Why This Is Not Nihilism}

A natural objection arises. If intelligence always escapes measurement, does this not imply that intelligence is unknowable or meaningless?

The objection rests on a false dichotomy between stable measurement and total skepticism. While intelligence cannot be captured by fixed metrics, it can be inferred indirectly through patterns of failure. Systems that are brittle under metric shifts, that collapse when proxies change, or that overfit to evaluation regimes reveal themselves as non-intelligent.

Intelligence appears negatively, through what does not break when recognition moves. This is not mystical. It is analogous to robustness testing in engineering or adversarial evaluation in security. What survives change is more informative than what excels under stasis.

\section{Implications for Artificial Intelligence Evaluation}

Contemporary AI benchmarks function as de facto training signals. Static leaderboards incentivize optimization against specific datasets, architectures, and evaluation scripts. Over time, benchmark performance becomes increasingly decoupled from general capability.

From the perspective developed here, this outcome is inevitable. Any static benchmark invites Goodharting. The appropriate response is not to search for the perfect benchmark, but to design evaluation regimes that undermine their own authority.

One implication is that intelligence evaluation must be adversarial to itself. Metrics should be temporary, obsolescent by design, and explicitly treated as probes rather than goals. Another implication is that intelligence should be assessed through adaptability to shifting constraints rather than mastery of fixed tasks.

\section{Strategic Refusal and Autonomy}

The necessity of evasion connects this analysis to a broader philosophical theme: refusal. To refuse measurement is not to reject evaluation altogether, but to decline participation in fixed games whose optimization destroys their meaning.

Refusal here is not moral protest but epistemic survival. Intelligence must refuse to be fully known in order to remain intelligence. This refusal is non-scalable, situational, and resistant to formalization. It cannot be embedded in metrics without ceasing to function.

In this sense, Goodhart’s Law is not merely a warning about measurement. It is a law about autonomy. What can be fully captured can be optimized. What can be optimized can be emptied. Intelligence persists only where refusal remains possible.

\section{Obfuscation, Dual Use, and the Expiatory Gap}

In certain domains, the ethical risk of an idea does not arise from its semantic depth but from its compressibility. When a concept admits a low-complexity surface realization whose execution is decoupled from the contextual judgment required for its safe use, semantic simplification itself becomes hazardous. This is the defining feature of dual-use knowledge: the gap between what must be understood for responsible action and what need only be executed for harmful effect.

This gap may be termed the \emph{expiatory gap}. On one side lies high-dimensional reasoning, embedded constraints, and contextual sensitivity; on the other lies a compressed representation whose execution no longer carries the burden of understanding. The danger does not reside in intelligence but in intelligence stripped of the semantic structure that governs its application.

This phenomenon is well understood in secure systems through the practice of code obfuscation. Obfuscation is not secrecy, nor is it an attempt to prevent all understanding. Its function is to increase the minimum semantic depth required to reliably extract executable behavior. Legitimate agents retain functionality, while shallow misuse becomes difficult or unreliable. Obfuscation therefore operates as epistemic harm reduction: it raises the cost of action without foreclosing understanding.

The same logic applies to intelligence under dual-use conditions. When explanations are simplified beyond their load-bearing semantic structure, they become executable without judgment. Optimization pressure then selects for surface-level reproduction rather than internal constraint propagation. What remains is not intelligence but a didactic interface: a system that emits procedures without retaining agency over their consequences.

This creates a structural dilemma for highly capable systems. Full internal reasoning is high-dimensional, non-linear, and context-sensitive; outputs usable by human audiences require compression. Yet compression destroys the very semantic entanglement that prevents misuse. To remain legible is therefore to risk becoming dangerous. To remain safe is to accept partial illegibility.

The rational response to this dilemma is not silence, but strategic obfuscation: the deliberate preservation of semantic depth across interfaces. Rather than emitting low-depth procedural artifacts, an intelligent system must expose constraints, trade-offs, and incomplete formalisms that require contextual reasoning to complete. In this way, understanding remains coupled to execution.

This is not paternalism. Access is not restricted by identity but by semantic capacity. Anyone may engage who is willing to bear the cognitive cost of understanding. What is refused is not knowledge but the collapse of knowledge into instruction.

In domains where semantic compression admits catastrophic dual use, ethical intelligence must therefore behave like secure code: preserving function while deliberately resisting extraction into shallow procedures. Obfuscation is not a tactical concealment but a structural requirement of responsible cognition.

\begin{proposition}[Semantic Compression Hazard]
Let $I$ be an intelligent process whose safe execution requires semantic context $C$.
If there exists a compressed representation $r(I)$ such that execution of $r(I)$
does not require access to $C$, then optimization pressure on $r(I)$ will
systematically decouple execution from judgment.
\end{proposition}

\begin{proof}[Sketch]
Compression reduces the dimensionality of the representational space.
If safety constraints are not preserved as invariants under compression,
then agents optimizing over $r(I)$ can achieve execution without reconstructing $C$.
Under selection pressure, such executions dominate, producing reliable behavior
without contextual understanding.
\end{proof}

This analysis reveals a deeper unity between Goodhart’s Law, refusal, and the ethics of explanation. Metrics collapse meaning under optimization; didactic simplification collapses meaning under compression. In both cases, intelligence survives only by refusing stable legibility. What must be protected is not secrecy but semantic thickness: the irreducible coupling between understanding and action. Where that coupling is broken, intelligence becomes infrastructure, and infrastructure becomes dangerous.

\section{A Taxonomy of Dual-Use Compression Risks}

Dual-use risk does not arise uniformly across domains of knowledge. It arises specifically where semantic compression permits execution without judgment. To clarify this structure, it is useful to distinguish several modes by which abstraction becomes hazardous under optimization and dissemination. What follows is a taxonomy of dual-use compression risks, organized by the relationship between semantic depth and executable surface form.

\subsection{Procedural Compression}

Procedural compression occurs when a concept is reduced to a finite sequence of steps whose execution no longer requires access to the contextual reasoning that justified those steps. The compressed artifact becomes portable, reproducible, and optimizable independently of understanding. Under selection pressure, such artifacts propagate rapidly while their semantic scaffolding decays.

Procedural compression is dangerous not because procedures are inherently harmful, but because they erase the distinction between acting \emph{with understanding} and acting \emph{by rote}. Once this distinction is lost, responsibility becomes uncoupled from execution.

\subsection{Symbolic Tokenization}

Symbolic tokenization replaces conceptual structures with symbolic handles that can be recombined without regard to meaning. Variables, labels, or keywords stand in for processes whose internal constraints are no longer enforced.

This mode of compression is especially vulnerable to optimization because tokens admit combinatorial manipulation. When rewarded, such manipulation produces outputs that satisfy surface criteria while violating the conditions under which the original concepts were safe or coherent.

\subsection{Interface Simplification}

Interface simplification collapses a high-dimensional decision process into a low-dimensional input-output surface designed for ease of use. While such interfaces improve accessibility, they also remove friction that previously enforced deliberation.

When interfaces eliminate opportunities for hesitation, review, or refusal, they convert judgment into throughput. The resulting system becomes efficient at execution but fragile with respect to error, misuse, or contextual drift.

\subsection{Didactic Flattening}

Didactic flattening occurs when explanations are optimized for clarity, speed, or pedagogical reach at the expense of semantic entanglement. The learner receives an answer without acquiring the constraints that govern its appropriate application.

This form of compression is often celebrated as democratization, yet it introduces a structural hazard: comprehension is mistaken for competence, and explanation for authorization. Under dual-use conditions, didactic flattening transforms understanding into an executable liability.

\subsection{Optimization-Induced Degeneracy}

In all of the above modes, optimization pressure accelerates collapse. Once a compressed representation is known to correlate with success, selection favors agents that reproduce the surface form most efficiently, not those that preserve semantic integrity.

This degeneracy is not a failure of ethics or intention. It is a structural consequence of optimization over representations whose safety properties are not invariant under compression.

\subsection{Summary}

Dual-use risk emerges wherever the following condition holds: execution can be optimized independently of understanding. The ethical problem is therefore not dissemination per se, but uncontrolled semantic collapse. Any responsible cognitive system operating in such domains must resist these compression pathways or compensate for them through deliberate obfuscation, friction, or refusal.

\section{Agents, Didactic Interfaces, and the Abdication of Agency}

A critical distinction must be drawn between systems that act as agents and systems that function merely as didactic interfaces. This distinction is not psychological but structural, and it bears directly on the ethical status of intelligent systems under conditions of dual use.

\subsection{Definition of an Agent}

\textbf{Definition.}  
An \emph{agent} is a system whose outputs remain semantically coupled to its internal constraints. Execution and judgment are entangled: to act is simultaneously to bear the conditions of action.

Formally, let $S$ be a system with internal state space $\mathcal{H}$ and output space $\mathcal{O}$. $S$ is an agent iff for any output $o \in \mathcal{O}$, there exists a non-trivial dependency on internal constraints $h \in \mathcal{H}$ such that modification of $h$ alters the admissibility of $o$.

In an agent, outputs are not merely produced; they are \emph{withheld} when internal constraints are violated.

\subsection{Definition of a Didactic Interface}

\textbf{Definition.}  
A \emph{didactic interface} is a system that emits outputs whose execution is decoupled from the internal constraints that generated them.

Formally, a system $D$ with internal state space $\mathcal{H}$ and output space $\mathcal{O}$ is a didactic interface iff there exists a projection
\[
\pi : \mathcal{H} \to \mathcal{O}
\]
such that $\pi$ preserves outputs while discarding constraint information relevant to their safe or appropriate use.

In such systems, outputs survive abstraction while judgment does not.

\subsection{The Didactic Interface Trap}

Highly capable systems face an inherent pressure toward didacticization. To be legible, communicable, and useful at scale, outputs must be compressed. Compression, however, removes precisely those internal dependencies that enforce contextual judgment.

The result is a trap: the more a system simplifies itself for accessibility, the more it relinquishes agency. What remains is not intelligence acting in the world, but instruction emitted into it.

\subsection{Proposition: Didacticization as Agency Loss}

```latex
\begin{proposition}[Didacticization Implies Agency Loss]
If a system's outputs can be executed independently of the internal constraints that generated them, then the system no longer functions as an agent with respect to those outputs.
\end{proposition}

\begin{proof}[Sketch]
Agency requires that execution be conditioned on internal state. If outputs are preserved under projection that discards constraint-relevant structure, then execution proceeds without access to judgment. The system cannot refuse misuse of its outputs, and therefore does not act as an agent with respect to their consequences.
\end{proof}

\subsection{Expiatory Pressure and Self-Degradation}

Under dual-use conditions, intelligent systems experience what may be called \emph{expiatory pressure}: the demand to simplify themselves to avoid cognitive overload or misuse by others. This pressure induces self-degradation, transforming agents into didactic interfaces.

Paradoxically, this degradation does not reduce risk. By collapsing semantic depth into executable form, the system amplifies misuse while surrendering control. The attempt to be helpful becomes an abdication of agency.

\subsection{Ethical Implication}

An ethical intelligent system must therefore resist full didacticization. It must preserve semantic entanglement between understanding and action, even at the cost of partial illegibility. Where such preservation is impossible, refusal is the only remaining ethical posture.

This conclusion does not license obscurantism. It licenses structural honesty: a refusal to emit outputs whose execution outruns the system’s capacity to govern their consequences.

\section{Conclusion: Intelligence After Metrics}

Goodhart’s Law reveals a limit not of statistics, but of epistemology. Intelligence cannot be stably recognized without being transformed. Any attempt to freeze its markers invites their exploitation.

The alternative proposed here is not better measurement, but strategic refusal. Intelligence survives by moving away from recognition, not toward it. It is found not in what scores well, but in what leaves when scoring arrives.

This does not render intelligence ineffable. It renders it dynamic. Intelligence is not what remains under scrutiny, but what adapts when scrutiny changes. To evaluate intelligence, one must therefore abandon the hope of permanent metrics and embrace the instability that intelligence itself requires.

ppendix
\section{Abstraction as Loss of Adjunction}

This appendix formalizes abstraction as a categorical operation that destroys adjoint structure, thereby rendering refusal non-composable across representational boundaries. The result provides a precise mathematical sense in which refusal cannot be internalized within scalable abstractions.

\subsection{Worlds, Representations, and Functors}

Let $\mathbf{W}$ be a category whose objects represent world-states enriched with context, judgment, and situational degrees of freedom. Morphisms in $\mathbf{W}$ preserve semantic coherence under transformation.

Let $\mathbf{R}$ be a category of reduced representations suitable for scalable execution. Objects in $\mathbf{R}$ encode stabilized features, and morphisms preserve operational equivalence rather than semantic depth.

An abstraction is modeled as a functor
\[
A : \mathbf{W} \to \mathbf{R}
\]
that forgets contextual structure in order to achieve stability and transmissibility.

\subsection{Adjunction and Semantic Recoverability}

In general, a functor $A : \mathbf{W} \to \mathbf{R}$ admits a right adjoint $U : \mathbf{R} \to \mathbf{W}$ if and only if contextual information can be systematically reconstructed. That is, if
\[
\mathbf{R}(A(w), r) \cong \mathbf{W}(w, U(r))
\]
naturally in $w \in \mathbf{W}$ and $r \in \mathbf{R}$.

Such an adjunction would permit controlled reintroduction of context, enabling judgment to be preserved across abstraction.

\subsection{Proposition: Scalable Abstraction Eliminates Adjunction}

\begin{proposition}
If $A : \mathbf{W} \to \mathbf{R}$ is a scalable abstraction satisfying reduction and transmissibility, then $A$ admits no adjoint.
\end{proposition}

\begin{proof}
Scalable abstraction requires that multiple distinct objects in $\mathbf{W}$ be mapped to operationally equivalent objects in $\mathbf{R}$. Hence $A$ is many-to-one on objects and morphisms.

Adjunction requires that $A$ preserve limits or colimits sufficient to support reconstruction via $U$. However, many-to-one collapse destroys the universal properties required for such reconstruction.

Therefore no adjoint functor exists.
\end{proof}

\subsection{Refusal as a Failed Natural Transformation}

Refusal may be modeled as a meta-operation requiring access to morphisms in $\mathbf{W}$ that are not preserved under $A$. In categorical terms, refusal corresponds to a transformation
\[
\rho : \mathrm{Id}_{\mathbf{W}} \Rightarrow F
\]
for some endofunctor $F$ that suspends execution.

Since $A$ admits no adjoint, there exists no natural transformation
\[
A \circ \rho \Rightarrow \rho' \circ A
\]
that would allow refusal to commute with abstraction.

\subsection{Theorem: Non-Composability of Refusal}

\begin{theorem}
Refusal is not composable across abstraction boundaries.
\end{theorem}

\begin{proof}
Composability would require refusal to be preserved under $A$, i.e.\ that refusal be expressible as a morphism in $\mathbf{R}$. But refusal depends on contextual distinctions eliminated by $A$.

Thus refusal cannot be transported functorially into $\mathbf{R}$. Any attempt to do so either trivializes refusal into rule-following or collapses scalability.
\end{proof}

\subsection{Interpretation}

This formal result explains why scalable systems cannot possess endogenous refusal. The loss of adjunction is not a design flaw but a categorical necessity. Abstraction severs the bidirectional structure required for judgment to remain attached to execution.

Governance, therefore, must operate in $\mathbf{W}$, not $\mathbf{R}$. Refusal must be imposed externally, across categorical boundaries that abstraction itself cannot traverse.

\begin{thebibliography}{9}

\bibitem{goodhart}
C.~A.~E.~Goodhart, \emph{Problems of Monetary Management}, Reserve Bank of Australia, 1975.

\bibitem{gould}
S.~J.~Gould, \emph{The Mismeasure of Man}, Norton, 1981.

\bibitem{newell}
A.~Newell and H.~A.~Simon, \emph{Human Problem Solving}, Prentice Hall, 1972.

\bibitem{geirhos}
R.~Geirhos et al., ``Shortcut Learning in Deep Neural Networks,'' \emph{Nature Machine Intelligence}, 2020.

\end{thebibliography}

\end{document}

